% =====================================================================
% Greybody factors for <SPACETIME-TBD> — working draft
%
% Build:    latexmk -pdf main.tex
% Watch:    latexmk -pvc -pdf main.tex
% Clean:    latexmk -c
%
% This is a *working draft*. The spacetime, field content, and numerical
% method are not yet locked in --- see ./README.md for the open decisions
% and the criteria for closing them.
% =====================================================================

\documentclass[
  aps,
  prd,
  preprint,      % single-column draft style; switch to 'reprint' near submission
  superscriptaddress,
  nofootinbib,
  amsmath,
  amssymb,
  floatfix,
]{revtex4-2}

% --- Packages -------------------------------------------------------
\usepackage{graphicx}
\usepackage{bm}
\usepackage{physics}
\usepackage{tensor}
\usepackage{xcolor}
\usepackage[colorlinks=true,
            linkcolor=blue!60!black,
            citecolor=green!60!black,
            urlcolor=blue!60!black]{hyperref}

\usepackage[
  backend=biber,
  style=numeric-comp,
  sorting=none,
  giveninits=true,
  uniquename=init,
]{biblatex}
\addbibresource{../../references/library.bib}

% --- Custom macros --------------------------------------------------
\newcommand{\tr}{\operatorname{tr}}
\newcommand{\Tr}{\operatorname{Tr}}
\newcommand{\Sch}{\text{Sch}}
\newcommand{\GN}{G_{\!N}}
\newcommand{\lp}{\ell_{\mathrm{P}}}
\newcommand{\Mp}{M_{\mathrm{P}}}

% Draft watermark — remove before submission
\usepackage{draftwatermark}
\SetWatermarkText{DRAFT}
\SetWatermarkScale{4}
\SetWatermarkLightness{0.92}

% =====================================================================
\begin{document}

\title{Greybody factors for \textbf{\textcolor{red}{[SPACETIME-TBD]}} black holes}

\author{Sameer Pant}
\email{spant@clearoneadvantage.com}
\affiliation{Affiliation TBD}

\date{\today}

\begin{abstract}
\textcolor{red}{[Draft abstract --- fill in after results are in.]}
We compute the greybody factors $\Gamma_{\omega\ell}(\omega)$ for a massless
scalar field on the \textcolor{red}{[SPACETIME-TBD]} background and extract the
asymptotic Hawking emission spectrum at infinity. We compare to the
Schwarzschild result of Ref.~\cite{Page1976} and discuss the implications of
the modified geometry on the low- and high-frequency tails of the emission.
\end{abstract}

\maketitle

% --- Introduction ---------------------------------------------------
\section{Introduction}
\label{sec:intro}

Hawking's discovery that a black hole of mass $M$ emits a thermal flux at
temperature
\begin{equation}
  T_H = \frac{\hbar c^3}{8\pi \GN M k_B}
  \label{eq:hawking-temp}
\end{equation}
established that black holes are genuine thermodynamic objects~\cite{Hawking1975,
BardeenCarterHawking1973}. The flux measured at infinity, however, is not
exactly Planckian: it is modulated by a frequency-dependent transmission
coefficient $\Gamma_{\omega\ell}(\omega)$ --- the \emph{greybody factor} ---
that encodes the back-scattering of outgoing modes by the curvature potential
outside the horizon~\cite{Page1976}.

\textcolor{red}{[Motivate the specific spacetime: why is this geometry
interesting? Where does it appear in the modern literature? What is known
about its Hawking emission already, and what is the gap?]}

The rest of the paper is organized as follows.
Section~\ref{sec:setup} introduces the geometry and the scalar field equation.
Section~\ref{sec:numerics} describes the numerical method for extracting
$\Gamma_{\omega\ell}(\omega)$.
Section~\ref{sec:results} presents the greybody factors and the emission
spectrum, and compares them to the Schwarzschild benchmark.
Section~\ref{sec:discussion} discusses the limits of the calculation and
points to extensions.

% --- Setup ----------------------------------------------------------
\section{Setup}
\label{sec:setup}

\subsection{The \textcolor{red}{[SPACETIME-TBD]} geometry}

\textcolor{red}{[Metric, parameters, horizon structure, asymptotic behavior.
Include the limit in which the geometry reduces to Schwarzschild.]}

\subsection{Scalar field equation and the radial potential}

The minimally-coupled massless Klein--Gordon equation
$\Box \phi = 0$, after separation of variables in
$\phi = e^{-i\omega t} Y_{\ell m}(\theta,\varphi)\, \psi_{\omega\ell}(r)/r$,
reduces to a Schr\"odinger-like equation
\begin{equation}
  \frac{d^2 \psi}{dr_*^2} + \left[\omega^2 - V_\ell(r)\right] \psi = 0,
  \label{eq:radial}
\end{equation}
with tortoise coordinate $r_*$ and effective potential $V_\ell$ \textcolor{red}{[
write the explicit form once the geometry is fixed]}.

\subsection{Boundary conditions}

The greybody factor $\Gamma_{\omega\ell}$ is defined by the standard scattering
boundary conditions: purely-ingoing at the horizon, mixed in/outgoing at
infinity. Concretely,
\begin{align}
  \psi &\sim e^{-i\omega r_*} & &(r \to r_+), \\
  \psi &\sim A^{\mathrm{out}}_{\omega\ell}\, e^{+i\omega r_*}
        + A^{\mathrm{in}}_{\omega\ell}\, e^{-i\omega r_*} & &(r \to \infty),
\end{align}
and
\begin{equation}
  \Gamma_{\omega\ell} = 1 - \left|\frac{A^{\mathrm{out}}_{\omega\ell}}
                                       {A^{\mathrm{in}}_{\omega\ell}}\right|^2.
  \label{eq:greybody-def}
\end{equation}

% --- Numerical method ------------------------------------------------
\section{Numerical method}
\label{sec:numerics}

\textcolor{red}{[Pick one of: (i) direct shooting, (ii) Leaver-style continued
fractions, (iii) WKB. Justify the choice. State accuracy targets and
convergence tests.]}

% --- Results ---------------------------------------------------------
\section{Results}
\label{sec:results}

\begin{figure}[t]
  \centering
  % \includegraphics[width=\columnwidth]{figures/greybody.pdf}
  \caption{\textcolor{red}{[Greybody factor $\Gamma_{\omega\ell}(\omega)$ for
  several $\ell$, with the Schwarzschild benchmark of Ref.~\cite{Page1976}
  overlaid.]}}
  \label{fig:greybody}
\end{figure}

\begin{figure}[t]
  \centering
  % \includegraphics[width=\columnwidth]{figures/spectrum.pdf}
  \caption{\textcolor{red}{[Hawking emission spectrum
  $d^2N/(dt\, d\omega)$ at infinity, summed over $\ell$.]}}
  \label{fig:spectrum}
\end{figure}

% --- Discussion ------------------------------------------------------
\section{Discussion}
\label{sec:discussion}

\textcolor{red}{[Implications for the emission spectrum's low/high-frequency
tails, comparison to Schwarzschild, what fails as the geometry approaches
limits (e.g. extremal, near-Schwarzschild). Punt anything that does not fit
in a short paper to "future work" --- additional fields, additional
spacetimes, QNM spectra, holographic interpretations all belong in a sequel.]}

% --- Acknowledgments -------------------------------------------------
\begin{acknowledgments}
\textcolor{red}{[TBD]}
\end{acknowledgments}

\printbibliography

% --- Appendices ------------------------------------------------------
\appendix

\section{Conventions}
\label{app:conventions}

Metric signature $(-,+,+,+)$. Natural units $\hbar = c = \GN = k_B = 1$ unless
stated. Greek indices $\mu,\nu,\dots$ run over spacetime; Latin indices run
over the angular sphere.

\section{Numerical convergence tests}
\label{app:convergence}

\textcolor{red}{[Tables / plots showing that the reported $\Gamma_{\omega\ell}$
are converged in resolution, integration domain, and numerical tolerance.]}

\end{document}
